\documentclass[10pt]{article}

% ------------------------------------------------------------
% COMP 2300 Proofs Core companion
% ------------------------------------------------------------
\usepackage[letterpaper,margin=0.58in,includehead,includefoot,
  headheight=15pt,headsep=8pt,footskip=22pt]{geometry}
\usepackage{amsmath,amssymb,mathtools}
\usepackage{microtype}
\usepackage[dvipsnames,svgnames]{xcolor}
\usepackage{multicol}
\usepackage{enumitem}
\usepackage{tabularx,booktabs,array}
\usepackage[most]{tcolorbox}
\usepackage{fancyhdr}
\usepackage{titlesec}
\usepackage{needspace}
\usepackage[hidelinks]{hyperref}

\definecolor{ProofNavy}{HTML}{17324D}
\definecolor{ProofTeal}{HTML}{176B6B}
\definecolor{ProofGold}{HTML}{A66B00}
\definecolor{ProofRed}{HTML}{9A3C35}
\definecolor{ProofBlueTint}{HTML}{EDF4F8}
\definecolor{ProofTealTint}{HTML}{EDF8F6}
\definecolor{ProofGoldTint}{HTML}{FFF7E5}
\definecolor{ProofRedTint}{HTML}{FCEFED}

\setlength{\parindent}{0pt}
\setlength{\parskip}{4pt}
\setlength{\columnsep}{0.28in}
\setlength{\columnseprule}{0.25pt}
\setlength{\emergencystretch}{1em}
\setlength{\abovedisplayskip}{5pt}
\setlength{\belowdisplayskip}{5pt}
\setlength{\abovedisplayshortskip}{3pt}
\setlength{\belowdisplayshortskip}{3pt}
\setlist{leftmargin=*,itemsep=2pt,topsep=3pt,parsep=0pt,partopsep=0pt}
\setlist[enumerate]{label=\arabic*.,ref=\arabic*}

\titleformat{\section}
  {\Large\bfseries\color{ProofNavy}}{}{0pt}{}
  [\vspace{-3pt}\color{ProofNavy}\rule{\linewidth}{0.7pt}]
\titleformat{\subsection}
  {\large\bfseries\color{ProofTeal}}{}{0pt}{}
\titleformat{\subsubsection}
  {\normalsize\bfseries\color{ProofNavy}}{}{0pt}{}
\titlespacing*{\section}{0pt}{5pt}{6pt}
\titlespacing*{\subsection}{0pt}{7pt}{3pt}
\titlespacing*{\subsubsection}{0pt}{5pt}{2pt}

\pagestyle{fancy}
\fancyhf{}
\fancyhead[L]{\textbf{COMP 2300}}
\fancyhead[C]{\textbf{Proofs Core Companion}}
\fancyhead[R]{Definitions, Rules, and Worked Examples}
\fancyfoot[L]{Core sections A--E}
\fancyfoot[C]{\thepage}
\fancyfoot[R]{Use definitions explicitly}
\renewcommand{\headrulewidth}{0.5pt}
\renewcommand{\footrulewidth}{0.3pt}

\tcbset{
  enhanced,
  before skip=5pt,
  after skip=5pt,
  boxrule=0.6pt,
  arc=1.5mm,
  left=6pt,right=6pt,top=5pt,bottom=5pt,
  fonttitle=\bfseries
}
\newtcolorbox{definitionbox}[1]{
  title={#1},colback=ProofBlueTint,colframe=ProofNavy,
  colbacktitle=ProofNavy,coltitle=white
}
\newtcolorbox{methodbox}[1]{
  title={#1},colback=ProofTealTint,colframe=ProofTeal,
  colbacktitle=ProofTeal,coltitle=white
}
\newtcolorbox{workedbox}[1]{
  title={#1},colback=white,colframe=ProofNavy,
  colbacktitle=ProofBlueTint,coltitle=ProofNavy
}
\newtcolorbox{warningbox}[1]{
  title={#1},colback=ProofRedTint,colframe=ProofRed,
  colbacktitle=ProofRed,coltitle=white
}
\newtcolorbox{checklistbox}[1]{
  title={#1},colback=ProofGoldTint,colframe=ProofGold,
  colbacktitle=ProofGold,coltitle=white
}

\newcommand{\term}[1]{\textbf{\color{ProofNavy}#1}}
\newcommand{\why}{\textit{Why this works.} }
\newcommand{\Z}{\mathbb Z}
\newcommand{\Q}{\mathbb Q}
\newcommand{\R}{\mathbb R}
\newcommand{\N}{\mathbb N}
\newcolumntype{Y}{>{\raggedright\arraybackslash}X}
\newcolumntype{M}{>{\raggedright\arraybackslash}p{0.29\linewidth}}

\begin{document}

\begin{center}
  {\Huge\bfseries\color{ProofNavy} Proofs Core Companion}\par
  \vspace{3pt}
  {\large Definitions, rules of inference, proof methods, and complete worked examples}\par
  \vspace{5pt}
  \color{ProofTeal}\rule{0.88\linewidth}{1.1pt}
\end{center}

\begin{checklistbox}{How to use this companion}
The Proofs Core is organized into five skills: \textbf{A} inference and proof
structure, \textbf{B} direct proofs, \textbf{C} contraposition and contradiction,
\textbf{D} cases, existence, and counterexamples, and \textbf{E} strategy choice
and proof repair. Study the vocabulary first, then cover a worked solution and
try to reproduce its logic without looking. On any proof, name the objects,
state their domains, unpack definitions, justify the central step, and finish
with the exact requested conclusion.
\end{checklistbox}

\section*{The Language of the Proofs Core}

\begin{multicols}{2}
\raggedcolumns
\subsection*{Statements, arguments, and logic}

\term{Statement (proposition).} A declarative sentence that is either true or
false. A \term{compound statement} is built from simpler statements with
logical connectives.

\term{Predicate.} A sentence with one or more input variables, such as
$P(x)$: ``$x$ is even.'' It becomes a statement after the variables are given
values or quantified.

\term{Domain (universe of discourse).} The collection of objects a variable is
allowed to represent. A \term{nonempty domain} contains at least one object.

\term{Variable; free and bound.} A variable is \emph{bound} when it lies within
the scope of a quantifier, as $x$ does in $\forall x\,P(x)$. Otherwise it is
\emph{free}. The \term{scope} of a quantifier is the expression it governs.

\term{Negation.} $\neg p$ means ``not $p$.'' It reverses the truth value of
$p$.

\term{Conjunction.} $p\land q$ means ``$p$ and $q$'' and is true only when both
parts are true.

\term{Disjunction.} $p\lor q$ means inclusive ``or'': at least one part is
true, possibly both.

\term{Conditional (implication).} $p\to q$ means ``if $p$, then $q$.'' Here
$p$ is the \term{hypothesis (antecedent)} and $q$ is the
\term{conclusion (consequent)}. It is false only when $p$ is true and $q$ is
false.

\term{Biconditional.} $p\leftrightarrow q$ means both $p\to q$ and $q\to p$;
it is read ``$p$ if and only if $q$.''

\term{Argument.} A list of \term{premises} offered as reasons for a
\term{conclusion}. The notation
\[
  P_1,P_2,\ldots,P_n\ \vdash\ C
\]
means that $C$ is derivable from the premises using valid rules.

\term{Valid argument.} An argument in which it is impossible for every premise
to be true while the conclusion is false. \term{Invalid} means such a situation
is possible. A \term{countermodel} specifies a domain/interpretation that makes
the premises true and the conclusion false.

\term{Rule of inference.} A truth-preserving argument form licensed for moving
from established statements to a new statement. A \term{derivation} is a
sequence of such justified steps.

\subsection*{Related conditionals}

For $p\to q$:
\begin{center}
\begin{tabularx}{\linewidth}{@{}lY@{}}
\toprule
\textbf{Name} & \textbf{Form} \\
\midrule
Original & $p\to q$ \\
Converse & $q\to p$ \\
Inverse & $\neg p\to\neg q$ \\
Contrapositive & $\neg q\to\neg p$ \\
\bottomrule
\end{tabularx}
\end{center}
Only the contrapositive is logically equivalent to the original conditional.

\columnbreak
\subsection*{Mathematical claims and proofs}

\term{Axiom (postulate).} A starting statement accepted without proof in the
mathematical system being used.

\term{Definition.} A precise agreement about the meaning of a term. In a
proof, a definition can be used in both directions.

\term{Theorem (proposition).} A mathematical statement proved from definitions,
axioms, and earlier results. A \term{lemma} is a supporting theorem; a
\term{corollary} follows quickly from another result. A \term{conjecture} is a
claim believed plausible but not yet proved.

\term{Proof.} A finite, logically valid argument that establishes a claim for
every object covered by the claim. An example can illustrate a universal claim,
but cannot prove it.

\term{Universal statement.} A claim of the form $\forall x\,P(x)$: every object
in the domain has property $P$.

\term{Existential statement.} A claim of the form $\exists x\,P(x)$: at least
one object in the domain has property $P$.

\term{Arbitrary object.} A representative object chosen with no special
property beyond the stated hypotheses. Proving the conclusion for it permits
universal generalization.

\term{Witness.} A particular value or object that makes an existential claim
true. A \term{fresh name} is a new symbol used temporarily for an existential
witness, without assuming which object it names.

\term{Direct proof.} To prove $p\to q$, assume $p$ and derive $q$.

\term{Proof by contraposition.} To prove $p\to q$, prove the equivalent
statement $\neg q\to\neg p$.

\term{Proof by contradiction.} Assume the negation of the desired conclusion
and derive an impossibility, such as $r\land\neg r$.

\term{Proof by cases.} Divide all possibilities into \term{exhaustive} cases
(cases that cover every allowed object), prove the claim in each case, and then
combine the cases.

\term{Existence proof.} A proof of $\exists x\,P(x)$. A
\term{constructive existence proof} displays a witness and verifies it. A
\term{nonconstructive existence proof} proves a witness must exist without
identifying one.

\term{Counterexample.} One allowed object for which a universal claim fails.
It disproves the claim $\forall x\,P(x)$ by establishing $\neg P(c)$ for a
specific $c$.

\term{Circular reasoning.} Treating the desired conclusion (or an equivalent
statement) as a premise. This does not prove the conclusion.

\term{Gap.} A missing or unjustified step. A \term{proof repair} identifies the
gap and supplies a valid replacement argument.

\term{Strategy.} The overall method chosen for a proof. A
\term{substantive step} is a mathematical move that begins the actual argument,
not merely ``use contradiction'' or ``let $n$ be an integer.''

\term{Restriction.} An added hypothesis that narrows the allowed inputs. The
\term{weakest necessary restriction} removes exactly the cases that make a
claim or proof fail, without excluding additional valid cases.
\end{multicols}

\clearpage
\section*{Definitions Used in the Proofs Core}

\begin{multicols}{2}
\raggedcolumns
\subsection*{Numbers and operations}

\begin{definitionbox}{Number systems}
\begin{itemize}
  \item $\N=\{0,1,2,\ldots\}$: the natural numbers in this course.
  \item $\Z=\{\ldots,-2,-1,0,1,2,\ldots\}$: the integers.
  \item $\Q$: the rational numbers.
  \item $\R$: the real numbers, consisting of rational and irrational numbers.
  \item \term{Positive} means $>0$; \term{nonnegative} means $\ge 0$.
\end{itemize}
\end{definitionbox}

\term{Sum, difference, product.} For numbers $a,b$, these are respectively
$a+b$, $a-b$, and $ab$.

\term{Consecutive integers.} Integers that occur one after another, such as
$n,n+1,n+2,n+3$.

\term{Largest and smallest elements.} An element $L$ is largest in a set $S$
if $L\in S$ and $x\le L$ for every $x\in S$. An element $s$ is smallest if
$s\in S$ and $s\le x$ for every $x\in S$.

\term{Even integer.} An integer $n$ for which
\[
 n=2k\quad\text{for some }k\in\Z.
\]
The integer $k$ is a witness to evenness.

\term{Odd integer.} An integer $n$ for which
\[
 n=2k+1\quad\text{for some }k\in\Z.
\]

\term{Divisibility.} For integers $a,b$ with $a\ne0$, $a$ \term{divides}
$b$, written $a\mid b$, when $b=ak$ for some $k\in\Z$. Then $a$ is a divisor
(factor) of $b$.

\term{Rational number.} A real number $r$ for which
\[
 r=\frac{a}{b}\quad\text{for some }a,b\in\Z\text{ with }b\ne0.
\]
Always state that the denominator is nonzero.

\term{Irrational number.} A real number that is not rational.

\term{Prime number.} An integer $p>1$ whose only positive divisors are $1$
and $p$.

\term{Composite number.} An integer $n>1$ that can be written $n=ab$ for
integers $a,b$ satisfying $1<a<n$ and $1<b<n$. Equivalently, it has a positive
divisor other than $1$ and itself.

\term{Perfect square.} An integer of the form $m^2$ for some integer $m$.

\term{Absolute value.} For $x\in\R$,
\[
 |x|=\begin{cases}
 x,&x\ge0,\\
 -x,&x<0.
 \end{cases}
\]
The phrase \term{with equality exactly when} means both directions must hold:
the equality holds if $x$ has the stated value, and only then.

\subsection*{Closure facts often used in proofs}
Integers are closed under addition, subtraction, and multiplication. Rational
numbers are closed under addition, subtraction, multiplication, and division by
a nonzero rational number. These facts justify claims such as ``$j-k\in\Z$''
or ``$(ad+bc)/(bd)\in\Q$.''

\columnbreak
\subsection*{Sets}

\term{Set.} A collection of objects. If $x$ is an object in $A$, write
$x\in A$; otherwise write $x\notin A$.

\term{Subset.} $A\subseteq B$ means every element of $A$ is also an element of
$B$:
\[
 A\subseteq B\quad\Longleftrightarrow\quad
 \forall x\,(x\in A\to x\in B).
\]
To prove an inclusion, begin with an arbitrary $x\in A$ and prove $x\in B$.

\term{Set equality.} $A=B$ means the sets have exactly the same elements.
Usually prove both $A\subseteq B$ and $B\subseteq A$.

\term{Element proof.} A proof of a set inclusion or equality that begins with
an arbitrary element, translates membership through set definitions, and ends
by placing that element in the target set.

\term{Intersection.} $A\cap B$ is the set of objects in both sets:
\[
 x\in A\cap B\quad\Longleftrightarrow\quad
 (x\in A\land x\in B).
\]

\term{Union.} $A\cup B$ is the set of objects in at least one set:
\[
 x\in A\cup B\quad\Longleftrightarrow\quad
 (x\in A\lor x\in B).
\]

\begin{workedbox}{Definition-writing model}
\textbf{Claim.} The sum of an even integer and an odd integer is odd.

\textbf{Solution.} Let $a$ be an arbitrary even integer and $b$ an arbitrary
odd integer. By definition, there are integers $j,k$ such that
$a=2j$ and $b=2k+1$. Then
\[
 a+b=2j+(2k+1)=2(j+k)+1.
\]
Because $j+k\in\Z$, the final expression has the defining form of an odd
integer. Therefore $a+b$ is odd.

\why Distinct arbitrary inputs receive distinct witnesses; the calculation
constructs the witness $j+k$ required by the definition of odd.
\end{workedbox}

\begin{warningbox}{Definitions are not decorative}
Do not write only ``$a$ and $b$ are even.'' Write $a=2j$ and $b=2k$ for
possibly different $j,k\in\Z$. Do not stop at $a+b=2j+2k$; factor it into
$2(j+k)$ and state that $j+k$ is an integer. The last sentence must match the
definition of the property being proved.
\end{warningbox}

\begin{checklistbox}{Definition checklist}
Before using a formula, ask:
\begin{enumerate}
  \item What is the domain of each object?
  \item Which definition introduces a witness?
  \item Are independently introduced witnesses allowed to be different?
  \item Is a denominator or divisor known to be nonzero?
  \item Does the final expression exactly match the target definition?
\end{enumerate}
\end{checklistbox}
\end{multicols}

\clearpage
\section*{A. Inference and Proof Structure}

\begin{definitionbox}{Reading an inference rule}
Expressions above the line are premises; the expression below is the licensed
conclusion. A rule is a reusable valid form. Logical equivalences, by contrast,
permit replacing an expression by an equivalent expression.
\end{definitionbox}

\begin{multicols}{2}
\raggedcolumns
\subsection*{The eight propositional rules}

\begin{center}
\renewcommand{\arraystretch}{1.28}
\begin{tabularx}{\linewidth}{@{}Y>{\centering\arraybackslash}p{0.42\linewidth}@{}}
\toprule
\textbf{Rule} & \textbf{Form} \\
\midrule
Modus Ponens (MP) & $p,\ p\to q\ \vdash\ q$ \\
Modus Tollens (MT) & $p\to q,\ \neg q\ \vdash\ \neg p$ \\
Hypothetical Syllogism (HS) & $p\to q,\ q\to r\ \vdash\ p\to r$ \\
Disjunctive Syllogism (DS) & $p\lor q,\ \neg p\ \vdash\ q$ \\
Addition (Add) & $p\ \vdash\ p\lor q$ \\
Simplification (Simp) & $p\land q\ \vdash\ p$ \\
Conjunction (Conj) & $p,\ q\ \vdash\ p\land q$ \\
Resolution (Res) & $p\lor q,\ \neg p\lor r\ \vdash\ q\lor r$ \\
\bottomrule
\end{tabularx}
\end{center}

For DS and Simplification, the symmetric version is also valid. For example,
$p\lor q$ and $\neg q$ yield $p$; $p\land q$ yields either conjunct.

\begin{workedbox}{Worked derivation with named rules}
\textbf{Premises:} $a\to b$, $b\to c$, $\neg c$. \quad
\textbf{Target:} $\neg a$.
\[
\begin{array}{rll}
1.&a\to b&\text{Premise}\\
2.&b\to c&\text{Premise}\\
3.&\neg c&\text{Premise}\\
4.&\neg b&\text{MT, lines 2 and 3}\\
5.&\neg a&\text{MT, lines 1 and 4}
\end{array}
\]
Every new line cites earlier lines and the rule that licenses it.
\end{workedbox}

\subsection*{Useful equivalences}
\[
\begin{aligned}
p\to q&\equiv \neg p\lor q,\\
p\to q&\equiv \neg q\to\neg p,\\
\neg(p\land q)&\equiv\neg p\lor\neg q,\\
\neg(p\lor q)&\equiv\neg p\land\neg q,\\
\neg\neg p&\equiv p.
\end{aligned}
\]
An equivalence is not itself one of the eight inference rules; label the law
used when rewriting.

\columnbreak
\subsection*{Invalid look-alikes}

\begin{warningbox}{Affirming the consequent}
\[
 p\to q,\ q\ \nvdash\ p
\]
Knowing the consequence occurred does not show that this particular hypothesis
caused it.
\end{warningbox}

\begin{warningbox}{Denying the antecedent}
\[
 p\to q,\ \neg p\ \nvdash\ \neg q
\]
The conclusion may still hold for some other reason.
\end{warningbox}

\begin{warningbox}{Illegal disjunction simplification}
\[
 p\lor q\ \nvdash\ p
\]
A conjunction guarantees each part; a disjunction guarantees only that at
least one part holds.
\end{warningbox}

\begin{workedbox}{Diagnosing validity}
``If a file is encrypted, it is unreadable without a key. This file is readable
without a key. Therefore it is not encrypted.''

Let $e$ mean ``encrypted'' and $u$ mean ``unreadable without a key.'' The form
is $e\to u$, $\neg u$, therefore $\neg e$: \textbf{modus tollens}, so the
argument is valid.

If the second premise were instead ``the file is unreadable,'' the form would
be $e\to u$, $u$, therefore $e$: affirming the consequent, so it would be
invalid.
\end{workedbox}

\subsection*{Universal shortcuts}
\term{Universal modus ponens:}
\[
\forall x(P(x)\to Q(x)),\quad P(a)\ \vdash\ Q(a).
\]
Use UI to obtain $P(a)\to Q(a)$, then MP.

\term{Universal modus tollens:}
\[
\forall x(P(x)\to Q(x)),\quad \neg Q(a)\ \vdash\ \neg P(a).
\]
Use UI, then MT.

\end{multicols}

\clearpage
\section*{Universal and Existential Instantiation Are Different}

\begin{checklistbox}{The central distinction}
A universal premise is a reusable guarantee about \emph{every} object, so it
may be instantiated at any chosen object. An existential premise promises only
that \emph{some} witness exists, so it receives a fresh temporary name; you may
not choose the witness or silently identify it with another existential
witness.
\end{checklistbox}

\begin{multicols}{2}
\raggedcolumns
\subsection*{The four quantified rules}

\begin{definitionbox}{Universal Instantiation (UI)}
\[
\forall x\,P(x)\ \vdash\ P(a)
\]
Here $a$ may be any particular object in the domain. A universal statement can
be reused for $a$, $b$, or any other object.
\end{definitionbox}

\begin{definitionbox}{Universal Generalization (UG)}
\[
P(c)\text{ proved for arbitrary }c\ \vdash\ \forall x\,P(x)
\]
The object $c$ must be genuinely arbitrary: it cannot depend on a special
example, an existential witness, or an extra assumption not shared by every
object in the domain.
\end{definitionbox}

\begin{definitionbox}{Existential Instantiation (EI)}
\[
\exists x\,P(x)\ \vdash\ P(c)\text{ for a fresh }c
\]
The name $c$ stands for one fixed but unknown witness. It is local to the
argument. Do not choose its value, reuse an existing name, or let the final
conclusion depend on which witness $c$ is.
\end{definitionbox}

\begin{definitionbox}{Existential Generalization (EG)}
\[
P(a)\ \vdash\ \exists x\,P(x)
\]
One verified instance is enough to establish existence. The object $a$ is the
witness.
\end{definitionbox}

\subsection*{Arbitrary object versus witness}
\begin{center}
\renewcommand{\arraystretch}{1.18}
\begin{tabularx}{\linewidth}{@{}MYY@{}}
\toprule
& \textbf{Arbitrary object} & \textbf{Existential witness} \\
\midrule
Source & Chosen for a universal proof & Promised by $\exists$ \\
Known? & Only general hypotheses & The existential property \\
Choose it? & Any object in domain & No; value is unknown \\
End move & UG may be possible & Remove fresh name from conclusion \\
\bottomrule
\end{tabularx}
\end{center}

\columnbreak
\subsection*{Complete valid derivation}

\begin{workedbox}{$\forall x(P(x)\to Q(x)),\ \exists xP(x)\ \vdash\ \exists xQ(x)$}
\[
\begin{array}{rll}
1.&\forall x(P(x)\to Q(x))&\text{Premise}\\
2.&\exists x\,P(x)&\text{Premise}\\
3.&P(c)&\text{EI from 2; fresh }c\\
4.&P(c)\to Q(c)&\text{UI from 1 at the same }c\\
5.&Q(c)&\text{MP from 3 and 4}\\
6.&\exists x\,Q(x)&\text{EG from 5}
\end{array}
\]
Line 4 uses the universal rule at the witness supplied on line 3. Line 6 removes
the temporary name from the conclusion.
\end{workedbox}

\subsection*{Why two existentials may have different witnesses}

\begin{warningbox}{Invalid inference}
\[
\exists xP(x),\quad \exists xQ(x)
\ \nvdash\ \exists x(P(x)\land Q(x)).
\]
The first EI must receive (say) fresh name $a$; the second must receive a
different fresh name $b$. Nothing says $a=b$.

\textbf{Countermodel.} Let the domain be $\{1,2\}$. Let $P$ hold only for $1$
and $Q$ hold only for $2$. Then both premises are true, but no object satisfies
both predicates, so the conclusion is false.
\end{warningbox}

\subsection*{Order of quantifiers changes meaning}

Over $\Z$,
\[
\forall x\,\exists y\,(y>x)
\]
is true: after an arbitrary $x$ is chosen, use witness $y=x+1$.
But
\[
\exists y\,\forall x\,(y>x)
\]
is false: it asks for one fixed integer larger than every integer. Given any
candidate $y$, the integer $y+1$ is larger.

\subsection*{Negating quantified statements}
\[
\neg\forall xP(x)\equiv\exists x\neg P(x),\qquad
\neg\exists xP(x)\equiv\forall x\neg P(x).
\]
Negation flips the quantifier and negates the predicate. For example, ``not
every integer is even'' means ``some integer is not even.''
\end{multicols}

\clearpage
\section*{B. Direct Proofs}

\begin{methodbox}{Direct-proof architecture for $\forall x\,(P(x)\to Q(x))$}
\begin{enumerate}
  \item Let $x$ be an arbitrary object in the stated domain and assume $P(x)$.
  \item Unpack $P(x)$ using definitions; introduce correctly scoped witnesses.
  \item Use algebra, definitions, axioms, and earlier theorems to build the
        defining form of $Q(x)$.
  \item Verify that any new coefficient, denominator, or witness lies in the
        required domain.
  \item State that $Q(x)$ follows. Because $x$ was arbitrary, the conditional
        holds for every allowed $x$.
\end{enumerate}
\end{methodbox}

\begin{multicols}{2}
\raggedcolumns
\subsection*{Worked parity proof}

\begin{workedbox}{The product of an even integer and any integer is even}
Let $a$ be an arbitrary even integer and let $b\in\Z$. Because $a$ is even,
$a=2k$ for some $k\in\Z$. Then
\[
 ab=(2k)b=2(kb).
\]
Because $k,b\in\Z$, closure of the integers under multiplication gives
$kb\in\Z$. Thus $ab$ has the form $2m$ for the integer $m=kb$, so $ab$ is
even.

\why The last two sentences explicitly identify and verify the witness required
by the definition of even.
\end{workedbox}

\subsection*{Worked rational-number proof}

\begin{workedbox}{The difference of rational numbers is rational}
Let $x,y\in\Q$. Then there are integers $a,b,c,d$ with $b\ne0$ and $d\ne0$
such that $x=a/b$ and $y=c/d$. Therefore
\[
 x-y=\frac{a}{b}-\frac{c}{d}=\frac{ad-bc}{bd}.
\]
The numerator $ad-bc$ and denominator $bd$ are integers. Since $b\ne0$ and
$d\ne0$, also $bd\ne0$. Hence the displayed fraction has the defining form of
a rational number, so $x-y\in\Q$.
\end{workedbox}

\columnbreak
\subsection*{Worked element proof}

\begin{workedbox}{$A\cap(B\cap C)\subseteq B$}
Let $x$ be an arbitrary element of $A\cap(B\cap C)$. By definition of
intersection,
\[
 x\in A\quad\text{and}\quad x\in B\cap C.
\]
Applying the intersection definition again to $x\in B\cap C$ gives $x\in B$
and $x\in C$. In particular, $x\in B$. Since every arbitrary element of
$A\cap(B\cap C)$ belongs to $B$, we conclude
$A\cap(B\cap C)\subseteq B$.

\why To prove a subset claim, start in the set on the left and finish in the
set on the right.
\end{workedbox}

\subsection*{Two-direction claims}
To prove $p\leftrightarrow q$, prove both $p\to q$ and $q\to p$. To prove
$A=B$, usually prove both inclusions. Never silently prove just one direction.

\begin{warningbox}{Frequent direct-proof errors}
\begin{itemize}
  \item Checking examples instead of using an arbitrary object.
  \item Reusing one witness for two independent inputs.
  \item Treating the target conclusion as an assumption.
  \item Producing $2k$ or $a/b$ without showing $k\in\Z$ or $b\ne0$.
  \item Beginning an inclusion proof with an element already in the target.
  \item Ending with algebra but not stating why it proves the requested property.
\end{itemize}
\end{warningbox}

\begin{checklistbox}{B-problem finishing sentence}
``Because \underline{\hspace{0.7in}} lies in the required domain, the final
expression has the defining form of \underline{\hspace{0.7in}}. Therefore
\underline{\hspace{0.8in}}.''
\end{checklistbox}
\end{multicols}

\clearpage
\section*{C. Contraposition and Contradiction}

\begin{multicols}{2}
\raggedcolumns
\subsection*{Proof by contraposition}

\begin{methodbox}{Template}
To prove $p\to q$:
\begin{enumerate}
  \item Write the contrapositive $\neg q\to\neg p$ explicitly.
  \item Assume $\neg q$.
  \item Use a direct argument to derive $\neg p$.
  \item State that the original implication follows because a conditional and
        its contrapositive are equivalent.
\end{enumerate}
Contraposition is useful when negating the conclusion produces a more workable
hypothesis. If that negation is a disjunction, handle both branches by cases.
\end{methodbox}

\begin{workedbox}{If $5n+1$ is even, then $n$ is odd}
We prove the contrapositive: \emph{if $n$ is even, then $5n+1$ is odd}.
Suppose $n$ is even. Then $n=2k$ for some $k\in\Z$. Hence
\[
 5n+1=5(2k)+1=10k+1=2(5k)+1.
\]
Because $5k\in\Z$, the number $5n+1$ is odd. Thus the contrapositive is true,
so the original implication is true.
\end{workedbox}

\subsection*{Negating compound conclusions}
If the conclusion is ``$m$ and $n$ are odd,'' its negation is ``$m$ is even or
$n$ is even'' (for integers). This disjunction naturally creates two cases.
Do not negate $q\land r$ as $\neg q\land\neg r$; De Morgan gives
$\neg q\lor\neg r$.

\columnbreak
\subsection*{Proof by contradiction}

\begin{methodbox}{Template}
To prove a statement $S$:
\begin{enumerate}
  \item Suppose, for contradiction, that $\neg S$.
  \item Use that supposition together with known facts.
  \item Derive an impossibility: $r\land\neg r$, an object both larger and not
        larger than another, or a rational number shown irrational.
  \item Conclude the supposition is false, so $S$ is true.
\end{enumerate}
For a conditional $p\to q$, a contradiction proof assumes $p\land\neg q$.
\end{methodbox}

\begin{workedbox}{There is no smallest positive rational number}
Suppose, for contradiction, that $s$ is the smallest positive rational number.
Then $s/2$ is rational because $s\in\Q$ and $2\ne0$. Also
\[
 0<\frac{s}{2}<s.
\]
Thus $s/2$ is a positive rational number smaller than $s$, contradicting the
choice of $s$. Therefore no smallest positive rational number exists.
\end{workedbox}

\begin{workedbox}{Rational plus irrational is irrational}
Let $r\in\Q$ and let $x\in\R\setminus\Q$. Suppose, for contradiction, that
$r+x\in\Q$. Rational numbers are closed under subtraction, so
\[
 x=(r+x)-r\in\Q.
\]
This contradicts the hypothesis that $x$ is irrational. Hence $r+x$ is
irrational.
\end{workedbox}

\begin{warningbox}{What is not a contradiction proof}
Repeating the conclusion, deriving something merely surprising, or writing
``contradiction'' without naming incompatible statements is insufficient.
State exactly which assumption conflicts with which established fact.
\end{warningbox}
\end{multicols}

\clearpage
\section*{D. Cases, Existence, and Counterexamples}

\begin{multicols}{2}
\raggedcolumns
\subsection*{Proof by exhaustive cases}

\begin{methodbox}{Template}
Identify cases whose disjunction covers every object. Prove the target in each
case under that case assumption. Then explicitly invoke exhaustiveness.
Common exhaustive splits include even/odd and $x\ge0$/$x<0$.
\end{methodbox}

\begin{workedbox}{$n(n+3)$ is even for every $n\in\Z$}
Every integer is even or odd.

\textbf{Case 1: $n$ is even.} Write $n=2k$. Then
$n(n+3)=2[k(n+3)]$, which is even because $k(n+3)\in\Z$.

\textbf{Case 2: $n$ is odd.} Write $n=2k+1$. Then
$n+3=2k+4=2(k+2)$, so $n(n+3)=2[n(k+2)]$, which is even.

The cases are exhaustive, so $n(n+3)$ is even for every integer $n$.
\end{workedbox}

\begin{workedbox}{$|x|\ge x$ for every real $x$}
\textbf{Case 1: $x\ge0$.} By definition, $|x|=x$, hence $|x|\ge x$.

\textbf{Case 2: $x<0$.} Then $|x|=-x>0>x$, so again $|x|\ge x$.

These cases cover every real number; therefore the claim holds for all
$x\in\R$.
\end{workedbox}

\columnbreak
\subsection*{Constructive existence}

\begin{methodbox}{Construct and verify}
For $\exists x\,P(x)$, provide a legal witness $c$ and check every part of
$P(c)$. For several objects, verify that they are distinct/consecutive and that
each has the required property.
\end{methodbox}

\begin{workedbox}{Three consecutive composite integers exist}
Use witnesses $48,49,50$. They are consecutive, and
\[
48=6\cdot8,\qquad49=7\cdot7,\qquad50=5\cdot10.
\]
Each factorization uses factors strictly between $1$ and the number, so every
witness is composite. Thus three consecutive composite positive integers exist.
\end{workedbox}

\subsection*{Counterexamples}

\begin{methodbox}{Disprove, do not merely doubt}
To disprove $\forall x\,P(x)$, choose one $c$ in the correct domain and verify
$\neg P(c)$. To disprove $p\to q$, make $p$ true and $q$ false.
\end{methodbox}

\begin{workedbox}{A difference of irrational numbers need not be irrational}
The universal claim ``the difference of two irrational numbers is irrational''
is false. Take $x=y=\sqrt2$. Both inputs are irrational, but
\[
 x-y=\sqrt2-\sqrt2=0,
\]
which is rational. This one legal input pair disproves the universal statement.
\end{workedbox}

\begin{warningbox}{Examples prove existence, not universality}
One successful example proves an existential claim. One failed example
disproves a universal claim. No finite list of successful examples proves a
universal claim unless an earlier theorem reduces the domain to those cases.
\end{warningbox}
\end{multicols}

\clearpage
\section*{E. Strategy Choice and Proof Repair}

\subsection*{Choose a strategy from the logical form}
\begin{center}
\renewcommand{\arraystretch}{1.25}
\begin{tabularx}{0.96\linewidth}{@{}p{0.18\linewidth}YY@{}}
\toprule
\textbf{Method} & \textbf{Promising signals} & \textbf{First substantive step} \\
\midrule
Direct & Definitions turn the hypothesis into useful algebra or membership facts. & Assume the hypothesis and unpack its definition. \\
Contraposition & Negating the conclusion gives a simpler defining form or useful cases. & State $\neg q\to\neg p$ exactly and assume $\neg q$. \\
Contradiction & Nonexistence, irrationality, uniqueness, or a negated claim gives a useful hypothetical object. & Assume the precise negation of the target. \\
Cases & Parity, signs, absolute value, or a disjunction naturally partitions the domain. & State an exhaustive case split. \\
Constructive existence & A concrete object can be built and checked. & Display the witness, then verify every property. \\
Counterexample & The claim is universal and appears false. & Choose a legal input that makes the conclusion fail. \\
\bottomrule
\end{tabularx}
\end{center}

\begin{multicols}{2}
\raggedcolumns
\subsection*{Repair 1: reused witnesses}

\begin{warningbox}{Flawed}
``Let $a,b$ be odd. Then $a=2k+1$ and $b=2k+1$, so $a-b=0$ is even.''

The same witness was imposed on independent arbitrary integers, falsely forcing
$a=b$.
\end{warningbox}

\begin{workedbox}{Repair}
Use possibly different $j,k\in\Z$: $a=2j+1$ and $b=2k+1$. Then
\[
 a-b=2j+1-(2k+1)=2(j-k).
\]
Since $j-k\in\Z$, the difference is even.
\end{workedbox}

\subsection*{Repair 2: division by zero}

\begin{warningbox}{Flawed}
``If $r$ is rational and $x$ irrational, then $rx$ is irrational, since
$x=(rx)/r$.'' The proof divides by $r$ without knowing $r\ne0$, and the claim
is false for $r=0$.
\end{warningbox}

\begin{workedbox}{Repair}
Correct the theorem: if $r$ is a \emph{nonzero} rational number and $x$ is
irrational, then $rx$ is irrational. If $rx$ were rational, then division by
the nonzero rational $r$ would give $x=(rx)/r\in\Q$, a contradiction.
\end{workedbox}

\columnbreak
\subsection*{Repair 3: circular reasoning}

\begin{warningbox}{Flawed}
To prove that the cube of an odd integer is odd, the writer begins ``Assume
$n^3$ is odd.'' This assumes the conclusion instead of the hypothesis.
\end{warningbox}

\begin{workedbox}{Repair}
Assume $n$ is odd, so $n=2k+1$ for some $k\in\Z$. Then
\[
n^3=(2k+1)^3
=2(4k^3+6k^2+3k)+1.
\]
The coefficient is an integer, so $n^3$ is odd.
\end{workedbox}

\subsection*{Repair 4: wrong arbitrary element}

\begin{warningbox}{Flawed}
To prove $A\subseteq B$ from $A\cap B=A$, starting with $x\in A\cap B$ only
proves something about elements already known to lie in $B$.
\end{warningbox}

\begin{workedbox}{Repair}
Let $x\in A$ be arbitrary. Since $A=A\cap B$, we have $x\in A\cap B$.
Therefore $x\in B$. Hence every element of $A$ lies in $B$, so $A\subseteq B$.
\end{workedbox}
\end{multicols}

\clearpage
\section*{One-Page Core Readiness Check}

\begin{checklistbox}{A fully correct Proofs Core response}
\small
Before finishing, verify that the response:
\begin{multicols}{3}
\begin{enumerate}
  \item states the hypotheses and begins with an arbitrary required object;
  \item identifies every domain and scopes variables/witnesses correctly;
  \item names inference rules when requested and justifies the central step;
  \item treats all cases, directions, or lettered parts requested;
  \item never divides without a nonzero condition or assumes the conclusion;
  \item reaches the exact requested conclusion in a final sentence.
\end{enumerate}
\end{multicols}
\end{checklistbox}

\begin{multicols}{2}
\raggedcolumns
\subsection*{Rule names at a glance}
\begin{center}
\renewcommand{\arraystretch}{1.18}
\begin{tabularx}{\linewidth}{@{}lY@{}}
\toprule
\textbf{Name} & \textbf{Recognize this move} \\
\midrule
MP & Use $p$ with $p\to q$ to obtain $q$. \\
MT & Use $\neg q$ with $p\to q$ to obtain $\neg p$. \\
HS & Chain $p\to q$ and $q\to r$. \\
DS & Eliminate one option from a disjunction. \\
Addition & Enlarge $p$ to $p\lor q$. \\
Simplification & Extract one conjunct from $p\land q$. \\
Conjunction & Combine established $p$ and $q$. \\
Resolution & Cancel complementary disjuncts. \\
\bottomrule
\end{tabularx}
\end{center}

\subsection*{Quantifier names at a glance}
\begin{center}
\renewcommand{\arraystretch}{1.18}
\begin{tabularx}{\linewidth}{@{}lY@{}}
\toprule
UI & Every object $\Rightarrow$ a chosen object. \\
UG & An arbitrary object $\Rightarrow$ every object. \\
EI & Some object $\Rightarrow$ a fresh temporary witness. \\
EG & One verified object $\Rightarrow$ some object exists. \\
\bottomrule
\end{tabularx}
\end{center}

\columnbreak
\subsection*{Definition forms at a glance}
\[
\begin{aligned}
n\text{ even}&\iff n=2k\text{ for some }k\in\Z,\\
n\text{ odd}&\iff n=2k+1\text{ for some }k\in\Z,\\
a\mid b&\iff b=ak\text{ for some }k\in\Z,\\
r\in\Q&\iff r=a/b,\ a,b\in\Z,\ b\ne0,\\
A\subseteq B&\iff \forall x(x\in A\to x\in B),\\
x\in A\cap B&\iff x\in A\land x\in B.
\end{aligned}
\]

\subsection*{Thirty-second error scan}
\begin{itemize}
  \item Did independent objects receive independent witnesses?
  \item Is an EI name fresh? Was a UG object genuinely arbitrary?
  \item Are all cases exhaustive and all requested directions complete?
  \item Is every division protected by a nonzero condition?
  \item Does a counterexample satisfy the hypothesis and fail the conclusion?
  \item Does the final sentence match the exact claim?
\end{itemize}
\end{multicols}

\begin{center}
\vfill
{\large\bfseries\color{ProofNavy} Proof writing is controlled use of definitions.}\par
{\small Choose objects carefully, preserve their scope, and make the final
sentence match the claim.}
\end{center}

\end{document}
